\blueprintfrontmatter

\begin{abstract}
\code{flt-regular-extended} certifies FLT for the \emph{regular} primes
$17 \le q < 350$ by a cheap route: a kernel Bernoulli-numerator check discharges
regularity, and \textsf{flt-regular} supplies FLT for regular primes. It reuses
the shared Bernoulli/Herbrand base of \code{flt-cyclotomic-nt} (external leaves).
\end{abstract}

\section{The Bernoulli regular route}

\begin{definition}[FLT for regular primes (external)]\label{ext:fltregular}
  \leanok
  \textsf{flt-regular}'s theorem: a regular prime satisfies FLT (via Kummer's
  ideal-descent argument).
\end{definition}

\begin{definition}[Bernoulli / Herbrand base (external)]\label{ext:bernbase}
  \leanok
  The kernel Bernoulli check \code{not\_irregular\_of\_bModN} and Herbrand's
  theorem, from \code{flt-cyclotomic-nt} --- the algebraic direction
  ``Bernoulli-regularity $\Rightarrow p \nmid h$''.
\end{definition}

\begin{theorem}[Regularity check $\Rightarrow$ FLT]\label{thm:regcheck}
  \lean{RegPrimes.flt_of_regCheck}
  \leanok
  \uses{ext:fltregular}
  A prime passing the regularity check satisfies \code{FermatLastTheoremFor}, the
  generic bridge instantiated per prime.
\end{theorem}

\begin{theorem}[Bernoulli-regularity $\Rightarrow$ FLT]\label{thm:bernoulli}
  \lean{FltOfBernoulli.fermatLastTheoremFor_of_bernoulli}
  \leanok
  \uses{thm:regcheck, ext:bernbase}
  If $p$ divides the numerator of none of $B_2,\dots,B_{p-3}$ (a kernel-\code{decide}d
  test), then $p$ is regular and FLT holds for $p$.
\end{theorem}

\begin{theorem}[First regular prime past 100]\label{thm:flt107}
  \lean{RegPrimes.flt_107}
  \leanok
  \uses{thm:bernoulli}
  \code{FermatLastTheoremFor}~$107$, the first regular exponent above $100$. All
  $47$ regular primes in $[17, 350)$ are proved by this same route (files
  \code{Flt17.lean} \dots\ \code{Flt349.lean}).
\end{theorem}
